\section{Análisis del código fuente de LlamaIndex: proceso de Indexing}

\subsection{Carga de documentos y Chunking}

\begin{enumerate}
    \item Se cargan todos los documentos del directorio Data (se obtiene una lista de Document)
    \item A cada documento se le aplica \textbf{Split}: el Chunk Size es de 1024 (en realidad alrededor de 997, restando los tokens de la ruta del documento), con un Overlap de 200
    \item El proceso de Split primero divide según sus propias reglas (produce alrededor de 759 splits) y luego los combina mediante \textbf{Merge} en los Chunk finales (alrededor de 22)
\end{enumerate}

\subsection{Creación de Node y Embedding}

\begin{enumerate}
    \item Cada Chunk crea un \textbf{Node} (con un node\_id único y el contenido de texto correspondiente)
    \item Se usa el Embedding Model (text-embedding-ada-002, con Batch Size 2048) para calcular el Embedding Vector de cada Node
    \item Al final se obtienen 22 Node, cada uno con un Embedding Vector de 1536 dimensiones ($512 \times 3$)
\end{enumerate}

\subsection{Resumen de la sección}

Proceso de Indexing: documento $\rightarrow$ Split $\rightarrow$ Merge $\rightarrow$ Chunk $\rightarrow$ Node $\rightarrow$ Embedding Vector. Cada paso tiene una implementación clara en el código fuente.
